% --- Page style for Answer Key ---
\pagestyle{fancy}
\fancyhf{}
\lhead{SAT Practice Exam -- Answer Key}
\rhead{}
\lfoot{}
\rfoot{Page \thepage}
\renewcommand{\headrulewidth}{0.4pt}
\renewcommand{\footrulewidth}{0.4pt}

\twocolumn[
  \begin{center}
    {\huge \textbf{SAT Practice Exam}} \\
    \vspace{0.2cm}
    {\Large \textbf{Answer Key}} \\
    \vspace{0.4cm}
  \end{center}
  \vspace{0.4cm}
]

% --- Part I: No Calculator ---
\section*{Part I: No Calculator SAT Math Questions}

\subsection*{Question 1}
$C = \frac{5}{9}(F - 32)$. Which of the following must be true? (I, II, III as in exam.)
\answerbox{(D) I and II only}

\textbf{Explanation:} Treat the equation as a line $C = \frac{5}{9}F - \frac{5}{9}(32)$. The slope is $\frac{5}{9}$, so an increase of 1 degree Fahrenheit gives an increase of $\frac{5}{9}$ degree Celsius (I is true). An increase of 1 degree Celsius corresponds to $\frac{9}{5} = 1.8$ degrees Fahrenheit (II is true). For III: an increase of $\frac{5}{9}$ degree Fahrenheit gives $C = \frac{5}{9}(\frac{5}{9}) = \frac{25}{81} \neq 1$, so III is false. The only choice with both I and II true is D.

\subsection*{Question 2}
$\displaystyle \frac{24x^2 + 25x - 47}{ax - 2} = -8x - 3 - \frac{53}{ax - 2}$, $x \neq \frac{2}{a}$. What is $a$?
\answerbox{(B) $-3$}

\textbf{Explanation:} Multiply both sides by $ax - 2$:
$24x^2 + 25x - 47 = (-8x - 3)(ax - 2) - 53$. Expand the right: $-8ax^2 - 3ax + 16x + 6 - 53 = -8ax^2 - 3ax + 16x - 47$. Equating $x^2$-coefficients: $-8a = 24$, so $a = -3$.

\subsection*{Question 3}
If $3x - y = 12$, what is the value of $\displaystyle \frac{8^x}{2^y}$?
\answerbox{(A) $2^{12}$}

\textbf{Explanation:} Write $8^x = (2^3)^x = 2^{3x}$. Then $\frac{8^x}{2^y} = \frac{2^{3x}}{2^y} = 2^{3x - y}$. Since $3x - y = 12$, the value is $2^{12}$.

\newpage
\subsection*{Question 4}
Points $A$ and $B$ on a circle of radius 1; arc $\stackrel{\frown}{AB}$ has length $\frac{\pi}{3}$. What fraction of the circumference is the arc?
\answerbox{$\dfrac{1}{6}$, or 0.166, or 0.167}

\textbf{Explanation:} Circumference $C = 2\pi r = 2\pi$. The fraction is $\dfrac{\pi/3}{2\pi} = \dfrac{1}{6}$.

\subsection*{Question 5}
$\displaystyle \frac{8 - i}{3 - 2i}$ in the form $a + bi$. What is $a$?
\answerbox{$2$}

\textbf{Explanation:} Multiply numerator and denominator by the conjugate $3 + 2i$. The numerator becomes $(8-i)(3+2i) = 24 + 16i - 3i - 2i^2 = 26 + 13i$ (since $i^2 = -1$); the denominator is $9 + 4 = 13$. So the quotient is $\frac{26+13i}{13} = 2 + i$. Hence $a = 2$.

\subsection*{Question 6}
Triangle $ABC$: $\angle B = 90^\circ$, $BC = 16$, $AC = 20$. Triangle $DEF$ is similar with each side $\frac{1}{3}$ of $ABC$. What is $\sin F$?
\answerbox{$\dfrac{3}{5}$ or 0.6}

\textbf{Explanation:} By the Pythagorean theorem, $AB = \sqrt{20^2 - 16^2} = \sqrt{144} = 12$. Since $DEF \sim ABC$ with $F \leftrightarrow C$, we have $\sin F = \sin C = \dfrac{\text{opp}}{\text{hyp}} = \dfrac{AB}{AC} = \dfrac{12}{20} = \dfrac{3}{5}$.

\newpage
\section*{Part II: Calculator-Allowed SAT Math Questions}

\subsection*{Question 7}
Left/right-handed table; 18 left-handed, 122 right-handed total. Probability a right-handed student selected at random is female?
\answerbox{(A) 0.410}

\textbf{Explanation:} Let $x$ = left-handed females, $y$ = left-handed males. Then right-handed females = $5x$, right-handed males = $9y$. So $x + y = 18$ and $5x + 9y = 122$. Solving: $x = 10$, $y = 8$. So $5(10) = 50$ right-handed females out of 122. Probability $= \frac{50}{122} \approx 0.410$.

\subsection*{Question 8}
Little's law: $N = rT$. Checkout: 84 shoppers/hour make a purchase, average 5 minutes in line. How many shoppers on average in the checkout line?
\answerbox{7}

\textbf{Explanation:} Rate $r = \frac{84}{60} = 1.4$ shoppers per minute; $T = 5$ minutes. So $N = rT = 1.4 \times 5 = 7$.

\subsection*{Question 9}
New store: 90 shoppers/hour enter, each stays 12 minutes. Average number in new store is what percent less than in original (45)?
\answerbox{60}

\textbf{Explanation:} New store: $r = \frac{90}{60} = 1.5$ per minute, $T = 12$, so $N = 1.5 \times 12 = 18$. Percent less $= \frac{45 - 18}{45} \times 100 = 60$.

\subsection*{Question 10}
$(p, r)$ on $y = x + b$; $(2p, 5r)$ on $y = 2x + b$. If $p \neq 0$, what is $\frac{r}{p}$?
\answerbox{(B) $\dfrac{3}{4}$}

\textbf{Explanation:} From $y = x + b$: $r = p + b$, so $b = r - p$. From $y = 2x + b$: $5r = 2(2p) + b = 4p + b$, so $b = 5r - 4p$. Equating: $r - p = 5r - 4p$ $\Rightarrow$ $3p = 4r$ $\Rightarrow$ $\frac{r}{p} = \frac{3}{4}$.

\newpage
\subsection*{Question 11}
Grain silo: cylinder (radius 5 ft, height 10 ft) and two cones (each radius 5 ft, height 5 ft). Closest volume in cubic feet?
\answerbox{(D) 1047.2}

\textbf{Explanation:} Cylinder: $V = \pi r^2 h = \pi (25)(10) = 250\pi$. Each cone: $V = \frac{1}{3}\pi r^2 h = \frac{1}{3}\pi(25)(5) = \frac{125\pi}{3}$. Total: $250\pi + 2 \cdot \frac{125\pi}{3} = \frac{1000\pi}{3} \approx 1047.2$ cubic feet.

\subsection*{Question 12}
$x = \text{avg}(m, 9)$, $y = \text{avg}(2m, 15)$, $z = \text{avg}(3m, 18)$. Average of $x$, $y$, $z$ in terms of $m$?
\answerbox{(B) $m + 7$}

\textbf{Explanation:} $x = \frac{m+9}{2}$, $y = \frac{2m+15}{2}$, $z = \frac{3m+18}{2}$. So $x + y + z = \frac{6m + 42}{2} = 3m + 21$. Average $= \frac{3m+21}{3} = m + 7$.

\subsection*{Question 13}
$f(x) = x^3 - x^2 - x - \frac{11}{4}$ graphed. If $f(x) = k$ has three real solutions, what could $k$ be?
\answerbox{$-3$}

\textbf{Explanation:} The horizontal line $y = k$ must cut the cubic in three places. From the graph, the only such horizontal line is at $y = -3$, so $k = -3$.

\subsection*{Question 14}
$q = \frac{1}{2}nv^2$. Ratio of dynamic pressure at velocity $1.5v$ to that at $v$?
\answerbox{2.25 or $\dfrac{9}{4}$}

\textbf{Explanation:} Let $q_1 = \frac{1}{2}n v^2$. Then $q_2 = \frac{1}{2}n (1.5v)^2 = 2.25 \cdot \frac{1}{2}n v^2 = 2.25\, q_1$, so the ratio $\frac{q_2}{q_1} = 2.25$.

\subsection*{Question 15}
$p(3) = -2$. Which must be true about $p(x)$?
\answerbox{(D) The remainder when $p(x)$ is divided by $x - 3$ is $-2$.}

\textbf{Explanation:} By the remainder theorem, when $p(x)$ is divided by $x - 3$, the remainder is $p(3) = -2$. So D must be true. Choices A, B, C are of the form ``$x - k$ is a factor,'' which would force $r = 0$ in $p(x) = (x - k)q(x) + r$; then $p(3) = (3-k)q(3)$ could equal $-2$ only for specific $q(3)$, not for all $p$ with $p(3) = -2$. Only D is always true.
